\section{The obstruction in odd cyclic covers}
\label{sec:formal}

The degree-two filling constructed in Section~\ref{sec:chambers} will
be contrasted with the following obstruction.  Throughout this section,
a marked filling has the boundary identification and the graph marking
specified there.

\begin{proposition}[Odd covers do not fill]
\label{for:odd-nonfilling}
For every odd integer $d\geq3$, the marked boundary $S_d$ admits no
marked filling.
\end{proposition}

Suppose, toward a contradiction, that such a filling exists.  We first
undo the law-curve surgeries and cut the resulting graph-type manifold
into pieces indexed by the plumbing graph.  For each cut, the peripheral
direction detected in degree-one cohomology determines the plus endpoint;
the other endpoint is minus.  We then construct the formal data of
Theorem~\ref{ten:assignment} using these designations.  They are independent
from edge to edge, so a vertex may meet both kinds of endpoint.

\subsection{Recovering the marked cut pieces}

Let $\Gamma_d$ denote the plumbing graph of the degree-$d$ cover.
Its vertices are indexed by $\mathbb Z/d$, and each consecutive pair
is joined by five red and five blue edges.  There are also $p$ named
extra loops at every vertex in the graph model of $X_d$; these loops
are not edges of $\Gamma_d$.  Since $d\geq3$, every plumbing edge has
two different endpoints, and the distinct ports of a law give distinct
incident edges.

\begin{lemma}[Marked cuts]
\label{for:cuts}
A marked filling of $S_d$ gives a compact connected topological
four-manifold $W$ with boundary identified with $\partial X_d$ and
with the marked graph homotopy type of $X_d$.  It admits a decomposition
over $\Gamma_d$ with connected vertex pieces $V_v$ and connected
three-dimensional cuts $Y_e$ such that:
\begin{enumerate}
\item the outside boundary face of $V_v$ is the prescribed $K_v$,
and $\partial Y_e$ is precisely the prescribed join torus;
\item every marked four-slot law at $v$ bounds a disk in $V_v$;
\item the cuts are compact smooth three-manifolds, and the filled
vertex pieces are smooth away from finitely many interior points.
\end{enumerate}
The identifications retain the based port meridians, their longitudes,
and the named extra loops.
\end{lemma}

\begin{proof}
We use the dual-handle and cut construction of \cite[Section~2.4]{PD},
keeping the full cyclic quotient instead of passing to its parity
quotient.  Attach to the proposed filling the dual two-handles that
undo surgery on the initial law curves.  The attaching circles
represent, with their prescribed paths, the lifted normal-circle
generators $z_j$.  Indeed $w(z_j)=0$, so every replacement piece and
each such generator has an individual lift.  These are free-factor
generators in the graph marking of the filling.  Attaching the
two-cells therefore cancels the corresponding one-cells in its
homotopy type, leaving the marked graph of $X_d$.

The boundary after these attachments is $\partial X_d$.  On the common
link exterior its map to the remaining graph is the original marking:
the quotient has killed exactly the normal-circle generators.  The
same holds on the restored solid tori, whose group elements come from
their gluing tori.  Equality of the resulting fundamental-group maps
is equality up to homotopy of graph-valued maps.  Thus the marking of
$W$ is the required one.  The two-handle cocores are disjoint proper
locally flat disks with the original law curves as boundaries.

For the cuts, pass to the cover associated to the quotient that forgets
the extra loops.  Its deck group is $\pi_1\Gamma_d$, acting freely on
the universal covering tree of $\Gamma_d$.  The fundamental group of
this cover of $W$ is freely generated by the named extra loops at tree
vertices, with their whiskers.  Prescribe an equivariant map to the
tree to be constant at the appropriate vertex on each outside face and
each cocore neighborhood, and to cross the appropriate edge on every
join-torus collar.  The prescriptions agree on overlaps.  They extend
over the remaining space because the associated tree bundle has
contractible fiber, as in the cut construction of \cite[Section~2.4]{PD}.

Choose the punctured smoothing used in that construction, with its
punctures away from the prescribed disks and collars.  Make the tree
map constant near each puncture.  Smooth approximation over the middle
of the edges and regular values then give compact smooth cuts,
disjoint from the disks.  It remains to obtain connected cuts and
connected vertex pieces; regular values alone do not ensure this.

Form the component graph of these cuts in the tree cover.  Its
fundamental group is a quotient of that of the cover: a graph path can
be lifted by joining successive crossings inside connected pieces.
Each named extra generator is conjugate to a loop lying in one piece,
and hence has trivial image in the component graph.  The component
graph is therefore a tree.  The
outside faces and prescribed torus collars embed the original plumbing
tree in it.  In fact different tree-vertex labels cannot lie in the
same component after the target midpoints have been removed, and the
prescribed collars give exactly the edges joining these principal
vertices.

Every component outside this principal subtree attaches at a unique
principal vertex.  A deck transformation stabilizing that component
fixes its attaching vertex, and hence is the identity.  There are only
finitely many orbits of cut and complementary components in the finite
quotient.  Since an orbit meets the branches attached to a fixed
principal vertex at most once, those branches have uniformly finitely
many edges and vertices.  Discard their closed cuts and merge their
pieces into the principal pieces.  This leaves exactly one connected
cut with each prescribed torus as its entire boundary.  The disk
prescriptions and boundary markings are unchanged.  Descending to
$\Gamma_d$ proves the lemma.
\end{proof}

The filled vertex pieces in this lemma need not be smooth.  We use
precisely the replacement for ordinary forms supplied by
\cite[Lemma~2.4]{PD}.  For completeness, if $Z_v$ is the finite puncture
set in $V_v$, take smooth functions on $V_v\setminus Z_v$ that are
separately constant near each puncture in degree zero, and smooth forms
vanishing near the punctures in positive degrees.  This is a
differential graded algebra computing $H^*(V_v;\mathbb C)$.  Viewed as
a sheaf on $V_v$, it has the constant sheaf as its degree-zero
cohomology sheaf and no higher cohomology sheaves.  Partitions of unity
can be made constant on smaller puncture neighborhoods, so the complex
is a fine resolution.  Requiring degree-zero functions to vanish at a
chosen point of $K_v$ removes only $H^0$.

All subsequent vertex recursions and gauges take place in this
algebra.  Each coefficient of a positive-degree formal form vanishes
near the punctures, so Stokes' formula is computed on a compact smooth
subdomain, with zero contribution from its additional end boundaries.
One common vanishing neighborhood for the entire formal series is
unnecessary: a fixed coefficient uses only finitely many terms.
Collar extensions and the coefficientwise divisions used below have
the same property.  Removing interior points in dimension four does
not change the fundamental group; gauges between flat systems may be
separately constant at these ends.  Thus the based holonomy calculation
uses the filled fundamental group and the filled cohomology.  Only the
unchanged compact smooth three-manifolds $Y_e$ enter the edge-potential
theorem.

\subsection{Restriction maps and peripheral directions}

All cohomology in this subsection has coefficients in $\mathbb C$.
Give the edges any temporary orientation, and let
\[
 \Delta^q:\bigoplus_v H^q(V_v)\longrightarrow\bigoplus_e H^q(Y_e)
\]
be the difference of restrictions at their two endpoints.

\begin{lemma}[Linear cut data]
\label{for:linear}
The map $\Delta^2$ is an isomorphism.  The map $\Delta^1$ is
surjective and restricts to an isomorphism on the subspace where
evaluation on every named extra loop is zero.  For each cut $Y_e$, the
image of $H^1(Y_e)$ in $H^1(\partial Y_e)$ is one of the two original
coordinate lines.
\end{lemma}

\begin{proof}
The graph homotopy type of $W$ gives $H^2(W)=H^3(W)=0$.
Mayer--Vietoris for the decomposition over $\Gamma_d$ consequently
makes $\Delta^2$ an isomorphism and $\Delta^1$ surjective.  Since the
pieces and cuts are connected, the kernel of $\Delta^1$ is
\[
 H^1(W)\big/H^1(\Gamma_d).
\]
The graph classes vanish on the named extra loops.  By the retained
marking, evaluation on those $dp$ loops has rank $dp$, and
$b_1(W)=b_1(\Gamma_d)+dp$.  These evaluations therefore identify the
displayed quotient with $\mathbb C^{dp}$.  Given any restrictions on
the cuts, its unique kernel correction makes all extra evaluations
zero, proving the second assertion.

For a connected oriented three-manifold with one torus boundary,
half-lives/half-dies makes the boundary image in degree one a line.
The restrictions from the two vertex ends span $H^1(Y_e)$.  At the
torus, each end restriction annihilates its longitude, which bounds
in its outside face $K_v$.  These two longitudes are the two different
coordinate slopes, because plumbing interchanges meridian and
longitude.  A line different from both coordinate lines meets each
of them trivially; if it were the boundary image, all vertex
restrictions would vanish at the boundary, contradicting surjectivity
onto that nonzero line.  Thus it is one coordinate line.
\end{proof}

For every $e$, denote by $\alpha_e$ the coordinate loop detected by
this line and by $\beta_e$ the other coordinate loop.  Call the endpoint
whose port meridian is $\alpha_e$ the \emph{plus endpoint}, denoted
$v_+(e)$; at the other, or \emph{minus endpoint} $v_-(e)$, the port
meridian is $\beta_e$.  At a plus endpoint $\beta_e$ bounds in $K_v$;
at a minus endpoint $\alpha_e$ bounds there.  Orient $Y_e$ as a face of
$V_{v_+(e)}$, and use the opposite orientation at the other end.  Put
\[
 E_+(v)=\{e:v_+(e)=v\},\qquad E_-(v)=\{e:v_-(e)=v\}.
\]
These designations are determined independently on the different
edges; in particular they impose no orientation on the underlying
cycle.

Set $\mathfrak g=\mathfrak{sl}_2(\mathbb C)$ with its invariant trace
pairing, and define the internal edge tangent space
\[
 \mathcal E_e=H^1(Y_e,\partial Y_e;\mathbb C)\otimes\mathfrak g,
 \qquad N_e=\dim_{\mathbb C}\mathcal E_e.
\]
Since both $Y_e$ and its boundary are connected, relative degree-one
cohomology identifies with the kernel of boundary restriction.  After
the extra evaluations have been fixed, evaluation on $\alpha_e$ comes
only from its plus endpoint.  Lemma~\ref{for:linear} therefore gives
\begin{equation}
\label{for:tangent-isomorphism}
 \bigoplus_v \mathcal A_v\xrightarrow{\ \cong\ }\bigoplus_e\mathcal E_e,
\end{equation}
where $\mathcal A_v$ is the subspace of $H^1(V_v)\otimes\mathfrak g$ with zero
extra evaluations and zero evaluations on its plus meridians; the map
is the signed difference of restrictions.  In particular all the
extra and plus evaluations are independent conditions on the vertex
spaces.  We shall also use the perfect pairing
\begin{equation}
\label{for:edge-pairing}
 H^1(Y_e,\partial Y_e)\otimes\mathfrak g
 \ \times\ H^2(Y_e)\otimes\mathfrak g\longrightarrow\mathbb C,
 \qquad (a,b)\longmapsto\int_{Y_e}\operatorname{tr}(a\wedge b).
\end{equation}
By Lemma~\ref{for:linear}, restriction of each vertex $H^2$ summand
to its incident cuts is injective.

\subsection{The formal data to be constructed}

For each vertex $v$, let $\mathcal L_v$ consist of the four-element
sets of incident edges containing two red and two blue edges.
The ports at $v$ correspond bijectively to incident edges, including
parallel edges, since $d\geq3$.  Order each such set into two red--blue
pairs.  Lemma~\ref{ch:laws} then supplies the independently conjugated
laws with leading bracket
$[[X_1,X_2],[X_3,X_4]]$, with a red and a blue input in each pair.
The following proposition is the interface with the tensor argument.
The ideals in its statement, and throughout its proof, are actual
ideals rather than radicals.

\begin{proposition}[Formal data from the cuts]
\label{for:formal-data}
For the cut pieces of Lemma~\ref{for:cuts}, there are a
characteristic-zero field $K$, coordinate blocks
$h_e=(h_{e,1},\ldots,h_{e,N_e})$, and algebraic power-series germs
$F_e(h_e),b_e(h_e)$ over $K$, with $F_e(0)=b_e(0)=0$, as follows.
For each vertex set
\begin{equation}
\label{for:vertex-potentials}
 \begin{split}
 T_v&=K[t_e:e\in E_+(v)]/(t_e^2:e\in E_+(v)),\\
 \mathcal S_v&=\sum_{e\in E_+(v)}(F_e+t_eb_e)
                  -\sum_{e\in E_-(v)}F_e.
 \end{split}
\end{equation}
There are smooth formal graph parametrizations $M_v(y_v,t)$ in the
incident $h_e$ blocks over $T_v$, centered modulo $(t)$, such that:
\begin{enumerate}
\item their central tangent maps are injective, and the difference map
from the sum of the vertex tangent spaces to the sum of the edge blocks
is an isomorphism;
\item $\mathcal S_v(M_v)=0$;
\item for $L\in\mathcal L_v$, with $I=L\cap E_+(v)$ and
$J=L\cap E_-(v)$,
\begin{equation}
\label{for:mixed-membership}
 \left(\prod_{i\in I}t_i\right)
 \left(\prod_{j\in J}b_j\right)(M_v)
 \in\bigl((\partial_h\mathcal S_v)(M_v)\bigr)
       \subset T_v[[y_v]].
\end{equation}
\end{enumerate}
Empty products equal one.  In every parameter ring, only the individual
squares $t_e^2$ vanish; mixed products are retained.  One can take
$K=\mathbb C((\sigma))$.
\end{proposition}

We prove this proposition in the next four subsections.  The edge
results are imported unchanged from \cite[Proposition~4.5 and
Theorem~5.1]{MT}.  What requires verification here is their assembly
at vertices with both kinds of endpoint, and the conversion of the
four-slot laws into \eqref{for:mixed-membership}.

We first use Stokes' formula to put the signed sum of edge potentials
at each vertex in the square of the ideal of its flatness equations.
A gradient calculation identifies this ideal with the restricted
gradient ideal and permits a correction to exact zero.  Next we pass
to algebraic central edge potentials and algebraic lifts of the marked
$\beta_e$-holonomies on their critical schemes, and recenter at a
common flat background.  The four-slot laws then put the required
products of plus parameters and logarithmic holonomy entries at minus
ends into that ideal.
Finally, we replace the logarithmic perturbations by polynomials in
those lifts and use a nilpotent correction to restore exact vanishing
while preserving those memberships.

\subsection{Flat backgrounds and the vertex identities}

On $\partial Y_e$ choose closed strip one-forms $\zeta_e,\xi_e$ with
periods respectively $(1,0)$ and $(0,1)$ on
$(\alpha_e,\beta_e)$, and put
$\epsilon_e=\int_{\partial Y_e}\zeta_e\wedge\xi_e$.
A trace in the $\xi_e$ polarization means a matrix multiple of
$\xi_e$.  Fix a basepoint on that torus and a nonzero
$T_e\in\mathfrak g$.  A \emph{based gauge} is a gauge transformation
equal to identity at this point.  For a formal family $s$ of connection
one-forms, its \emph{curvature ideal} is generated by the scalar
coefficient functions of $ds+\tfrac12[s,s]$.  The input below gives a
finite residual presentation of this ideal and identifies it with a
gradient ideal.  The formal critical scheme of a potential means the
formal germ cut out by that actual gradient ideal.

\begin{lemma}[Single-cut input from MT]
\label{for:edge-input}
Let $Y_e$ be compact, connected, oriented and smooth, with connected
torus boundary and based coordinate slopes $\alpha_e,\beta_e$.
Suppose the image of $H^1(Y_e;\mathbb C)$ on the boundary detects
$\alpha_e$ and vanishes on $\beta_e$.  With internal coordinates
$h_e$ on $\mathcal E_e$ and a scalar parameter $\tau^2=0$, there is
a formal family $s_e(h_e,\tau)$ of connection one-forms, centered at
the trivial connection, and formal functions
$f_e\in(h_e)^3$, $b_e^{\mathrm{form}}\in(h_e)^2$ such that:
\begin{enumerate}
\item The boundary trace is
$\tau T_e\zeta_e+v_e(h_e,\tau)\xi_e$.  The internal gradient ideal
of $f_e+\tau b_e^{\mathrm{form}}$ equals the full curvature ideal of
$s_e(h_e,\tau)$ in $\mathbb C[[h_e]][\tau]/(\tau^2)$.
\item Its quotient represents, over every local Artinian coefficient
algebra, flat connections modulo based gauges with
$G_{\alpha_e}=\exp(-\tau T_e)$; equivalently, it is the formal germ
of based representations of $\pi_1Y_e$ with that condition.
On this quotient $G_{\beta_e}=\exp(-v_e(h_e,\tau))$.
Once a flat connection has the displayed torus trace, its gauge into
this family is equal to identity on the whole torus.
\item At $\tau=0$, writing $G_{\beta_e}$ for the marked holonomy on
the central critical scheme, one has
\begin{equation}
\label{for:edge-log-coefficient}
 b_e^{\mathrm{form}}
 =-\epsilon_e\operatorname{tr}(T_e\log G_{\beta_e})
              \pmod{\operatorname{Jac}(f_e)}.
\end{equation}
\end{enumerate}
\end{lemma}

This is \cite[Proposition~4.5, with the based-gauge identification of
Lemma~4.4]{MT}.  In its convention the connection is $d+s_e$, accounting
for the minus signs in the holonomies.  We call the family $s_e$ a
\emph{slice}; the assertion about the quotient says that it represents
all such based gauge classes with the prescribed peripheral holonomy.
The hypotheses have been verified for our
cuts in Lemmas~\ref{for:cuts} and~\ref{for:linear}.  We use this input
with $\tau=t_e$ at a plus end and with $\tau=0$ at a minus end, always
in the same central edge coordinates.

We first work jointly formally at a parameter $\sigma=0$, before
inverting $\sigma$.  Choose the standard basis
$(e_0,h_0,f_0)$ of $\mathfrak{sl}_2(\mathbb C)$, with
$[e_0,f_0]=h_0$, $[h_0,e_0]=2e_0$, and $[h_0,f_0]=-2f_0$, and put
\[
 g_{ab}(\sigma)=\exp(a\sigma e_0)\exp(b\sigma f_0),
                  \qquad a,b\in\{0,1,2\}.
\]
By \cite[Lemma~6.5]{MT}, the nine operators
$\operatorname{Ad}(g_{ab})$ span
$\operatorname{End}_{\mathbb C((\sigma))}(\mathfrak g\otimes
\mathbb C((\sigma)))$.  Assign the eight nonidentity matrices to
eight named extra loops at each vertex, and assign identity to the
remaining extra loops.  All port-meridian holonomies are initially
trivial.  These assignments extend to a common flat background on
the marked ambient graph: the extra loops are independent vertex
generators in its free groupoid, and the transports along plumbing
edges can be chosen freely.  Choose them so the frames on the two
sides of each join agree with the prescribed basing paths.  A cycle
in $\Gamma_d$ contributes a free groupoid generator, not an additional
relation among these choices.
For the later law tests, enumerate these nine matrices as
$g_1,\ldots,g_9$.

The vertex construction of \cite[Section~6.2]{MT} applies to the
based form complex described above.  Choose a cohomological
contraction $(\iota_v,\pi_v,\mathsf K_v)$ there, with degree-zero
cohomology removed: $\iota_v$ chooses cohomology representatives,
$\pi_v$ projects onto cohomology, and $\mathsf K_v$ has degree $-1$,
with
\[
 d\mathsf K_v+\mathsf K_vd=1-\iota_v\pi_v,
 \quad \pi_v\iota_v=1,
 \quad \mathsf K_v^2=\mathsf K_v\iota_v=\pi_v\mathsf K_v=0.
\]
Its recursion in the full degree-one coordinates
$y_v^{\mathrm{full}}$ is
\[
 a_v=\iota_vy_v^{\mathrm{full}}
             -\tfrac12\mathsf K_v[a_v,a_v],\qquad
 \mathcal F_v=da_v+\tfrac12[a_v,a_v],\qquad
 r_v=\pi_v\mathcal F_v.
\]
The finite list $r_v$ has
$q_v=\dim(H^2(V_v;\mathbb C)\otimes\mathfrak g)$ entries.  Bianchi's
identity and the slice equation give
\begin{equation}
\label{for:curvature-list}
 \mathcal F_v=(1+\mathsf K_v\operatorname{ad}_{a_v})^{-1}\iota_v r_v.
\end{equation}
Thus the curvature ideal is generated by this list, whose entries have
order at least two.  The leading coefficient forms against the list
are the chosen representatives of $H^2(V_v)\otimes\mathfrak g$.
No independence or regular-sequence assertion about its entries is
needed.

Prescribe the chosen logarithmic holonomies on the extra loops, and
prescribe $-t_eT_e$ on every plus meridian at $v$.  The linearization
of based logarithmic holonomy is signed evaluation of a one-form.
Equation~\eqref{for:tangent-isomorphism} consequently makes these
conditions a smooth formal parameter slice.  Write its remaining
coordinates as $y_v$, and set
\[
 R_v=\mathbb C[[\sigma,y_v]][t_e:e\in E_+(v)]/(t_e^2),
 \qquad \mathcal I_v=(r_v)\subset R_v.
\]
The list has been substituted into this ring.  Its entries remain in
the square of the joint maximal ideal
$\mathfrak m_v=(\sigma,y_v,t)$.
Modulo $\mathcal I_v$ the vertex system is flat.  At a plus end its
peripheral holonomies are the prescribed $\alpha_e$ holonomy and
trivial $\beta_e$ holonomy; at a minus end the $\alpha_e$ holonomy is
trivial and the $\beta_e$ holonomy varies.  Therefore the minus
restriction uses the \emph{central} edge slice even if other ends of
that vertex carry parameters.

Let
\[
 \mathcal S_v^{\mathrm{form}}
 =\sum_{e\in E_+(v)}(f_e+t_eb_e^{\mathrm{form}})
                         -\sum_{e\in E_-(v)}f_e.
\]
We need a smooth section in the edge coordinates, not merely their
values on the residual quotient $R_v/\mathcal I_v$.  The following
is the local extension of \cite[Lemma~6.3]{MT}.

\begin{lemma}[Vertex identities with both endpoint signs]
\label{for:vertex-identities}
The flat edge restrictions lift to a formal parametrization
$M_v^{\mathrm{form}}$ over $R_v$ satisfying
\begin{equation}
\label{for:square-gradient}
 \mathcal S_v^{\mathrm{form}}(M_v^{\mathrm{form}})\in\mathcal I_v^2,
 \qquad
 (\partial_h\mathcal S_v^{\mathrm{form}})(M_v^{\mathrm{form}})
                         =L_vr_v,
\end{equation}
where $L_v$ has a full-column minor that is a unit at the joint origin.
With $\sigma$ and the $t_e$ fixed, the tangent maps at the joint origin
are the restriction maps in \eqref{for:tangent-isomorphism}.
\end{lemma}

\begin{proof}
We give the boundary calculation to distinguish the present statement
from the two-vertex conclusion of \cite[Proposition~6.1]{MT}.
For a Lie-algebra-valued one-form $a$ on an oriented three-manifold
$Y$, the action and its coefficient variation are
\begin{equation}
\label{for:action-convention}
 \begin{split}
 \mathcal C_Y(a)&=\int_Y\operatorname{tr}
       \left(\tfrac12a\wedge da+\tfrac16a\wedge[a,a]\right),\\
 \dot{\mathcal C}_Y(a)&=\int_Y\operatorname{tr}
       (\dot a\wedge\mathcal F_a)
       -\tfrac12\int_{\partial Y}\operatorname{tr}(a\wedge\dot a),
 \qquad \mathcal F_a=da+\tfrac12[a,a].
 \end{split}
\end{equation}
For the edge slice $s_e$ with trace $u_e\zeta_e+v_e\xi_e$ and
$u_e=t_eT_e$, the adjusted potential of
\cite[Proposition~4.5]{MT} is precisely
\[
 f_e+t_eb_e^{\mathrm{form}}
   =\mathcal C_{Y_e}(s_e)
          +\tfrac{\epsilon_e}{2}\operatorname{tr}(u_ev_e).
\]
This normalization will cancel the fixed-$u_e$ boundary variation.

On the free outside face $K_v$, normalize the flat quotient connection to
its strip representative: a sum of the signed based holonomy logarithms
times closed one-forms dual to the free generators, with disjoint supports.  The
disjoint meridian-dual disks and extra-loop-dual spheres in the standard
outside face supply these strips.  Each boundary torus meets only its
own port strip.  Thus the choices at different tori coexist in one
representative: a plus port supplies a $\zeta_e$ trace and a minus
port a $\xi_e$ trace.  This representative has zero Chern--Simons
action.  Its torus trace is
$u_e\zeta_e$ at a plus end, where $u_e=t_eT_e$, and is in the
$\xi_e$ polarization at a minus end.  The strip forms and torus
normalizations are those of \cite[Sections~4 and~6.3]{MT}.

After this normalization, gauges equal to identity on the tori put
the flat restrictions on the cuts into their central or perturbed edge
slices.  The boundary traces agree exactly in the residual quotient:
at a plus end $G_{\beta_e}=1$ forces $v_e=0$ there, while at a minus
end both traces have zero $\alpha_e$ coefficient and the same based
$\beta_e$ logarithm.  The torus-identity slice gauge of
\cite[Lemma~4.4]{MT} therefore applies at every end, and the gauges on
intersecting faces agree.  Lift the edge
coordinates and gauge logarithms from $R_v/\mathcal I_v$ to $R_v$,
and extend the logarithms through collars.  The lifting and division
argument in \cite[Section~6.3]{MT} uses a fixed finite residual list:
coefficientwise formal division expresses each discrepancy as an
actual sum $\sum_a r_{v,a}\gamma_a$ of form-valued series.  Every
output coefficient uses finitely many input coefficients.  This also
works in the punctured form complex.  Finally change the lifted
connection by such ideal multiples so that its restriction to $K_v$
is exactly the strip representative.  Call the result
$\widehat a_v$.

The curvature of $\widehat a_v$ lies in $\mathcal I_v$, and its edge
restrictions agree with the chosen slices modulo that ideal.  The
added ideal multiples change the leading curvature coefficient forms
in \eqref{for:curvature-list} only by exact forms; the gauges have
identity constant term.  Stokes' formula gives
\[
 \mathcal C_{\partial V_v}(\widehat a_v)
   =\tfrac12\int_{V_v}\operatorname{tr}
                 (\mathcal F_{\widehat a_v}\wedge
                  \mathcal F_{\widehat a_v})\in\mathcal I_v^2.
\]
The outside-face contribution is zero.  Interpolating an edge
restriction to its chosen slice has an interior action variation in
$\mathcal I_v^2$: the displacement and the curvature are both in
$\mathcal I_v$.  At a minus end the boundary variation is zero because
both traces are in the $\xi_e$ polarization.  At a plus end the traces
change from $u_e\zeta_e$ to $u_e\zeta_e+v_e\xi_e$, and the torus
variation is
\[
 -\tfrac{\epsilon_e}{2}\operatorname{tr}(u_ev_e).
\]
The adjustment in the perturbed edge potential cancels this term
exactly.  Reversing the orientation of a boundary face changes its
action by a sign.  Summing these per-end formulas proves the first
identity of \eqref{for:square-gradient}, including all mixed parameter
terms.

The internal edge variation formula pairs slice curvature with a
derivative of the slice.  It gives the second identity with an actual
coefficient matrix against the fixed list $r_v$.  At the joint origin
its entry for an edge tangent representative $\theta_{e,k}$ and a
vertex residual representative $\omega_{v,a}$ is, with the face sign,
\[
 \int_{Y_e}\operatorname{tr}
           (\theta_{e,k}\wedge\omega_{v,a}|_{Y_e}).
\]
Here $\theta_{e,k}$ is a closed relative one-form.  Added exact
two-forms do not change this pairing, by integration by parts and its
zero boundary trace.  Injectivity of the vertex $H^2$ restriction and
the perfect pairing \eqref{for:edge-pairing} make this matrix have full
column rank.  Its selected minor is a unit.  This argument computes
the coefficient matrix from curvature, and so does not assume that
the residual list has no syzygies.  Finally, every lift has the same
first derivatives as its residual-quotient restriction because
$\mathcal I_v$ has order at least two.  These are exactly the
cohomological tangent maps already identified.
\end{proof}

We can now arrange exact zero action without changing a flat marking.
Apply \cite[Lemma~6.4]{MT} separately at each vertex to
\eqref{for:square-gradient}.  Its hypotheses are visible here:
$\mathcal I_v$ has order at least two, and the internal Hessian of
$\mathcal S_v^{\mathrm{form}}$ vanishes at the joint origin, since the
central potentials have order at least three and all other terms carry
a parameter.  The correction is by $\mathcal I_v$-multiples, preserves
the tangent, and leaves the gradient ideal equal to $\mathcal I_v$.

More explicitly, a unit full-column minor of $L_v$ expresses an error
in $\mathfrak m_v^n\mathcal I_v^2$ as the internal gradient paired
with a displacement in $\mathfrak m_v^n\mathcal I_v$.  Cancel that
linear error.  The Hessian condition puts the Taylor remainder in
$\mathfrak m_v^{2n+1}\mathcal I_v^2$ and changes the gradient
coefficient matrix only by a matrix in
$\mathfrak m_v^{n+1}$.  Iteration converges and preserves its unit
minor.  Holding the list $r_v$ fixed throughout gives corrected
sections with
\begin{equation}
\label{for:exact-before-algebraic}
 \mathcal S_v^{\mathrm{form}}(M_v^{\mathrm{form}})=0,
 \qquad
 \bigl((\partial_h\mathcal S_v^{\mathrm{form}})
                  (M_v^{\mathrm{form}})\bigr)=\mathcal I_v.
\end{equation}
When $q_v=0$, the ideal is zero and no correction is needed.
Everything so far is local to a vertex except for the common edge
coordinates and background.  In particular it applies to a vertex
with no plus ends as well.
From now on write $M_v$ for these corrected sections, retaining that
notation when the sections are transported through subsequent
coordinate changes and recentering.

\subsection{Algebraic coordinates and recentering}

We have exact formal equations and their actual gradient ideals.  The
tensor theorem also requires algebraic edge functions.  We use the
marked algebraicity theorem \cite[Theorem~5.1]{MT}, whose formal
comparison is \cite[Proposition~5.2]{MT}.  Separately on every edge,
it gives an invertible formal coordinate change making the central
potential an algebraic germ $F_e$, together with algebraic matrix
lifts $G_e$ of the marked $\beta_e$ holonomy on its full critical
scheme.  Choose $G_e(0)=1$.  Transport all vertex sections and residual
identifications along these changes.

In these coordinates, \eqref{for:edge-log-coefficient} says that the
perturbation coefficient differs from
\[
 \ell_e(G_e)=-\epsilon_e\operatorname{tr}(T_e\log G_e)
\]
by an element of $\operatorname{Jac}(F_e)$.  A separate internal
coordinate change linear in $t_e$ absorbs that difference exactly.
Write $b_e^{\mathrm{form}}-\ell_e(G_e)
=\sum_k c_{e,k}\partial_{h_{e,k}}F_e$, and introduce new coordinates
$x_e=h_e+t_ec_e(h_e)$.  The identity
\[
 F_e(h_e+t_ec_e(h_e))
 =F_e(h_e)+t_e\sum_kc_{e,k}(h_e)\partial_{h_{e,k}}F_e(h_e)
\]
is exact because $t_e^2=0$, so the potential expressed in $x_e$ is
$F_e(x_e)+t_e\ell_e(G_e(x_e))$.  Rename $x_e$ as $h_e$.
Such a change is made only at the plus
end carrying that parameter, and is identity at parameter zero.
Thus both endpoints retain the same central coordinates and potential;
the minus end still uses just $-F_e$.  Exact vanishing, the gradient
ideal, and the quotient marking are transported with the sections.

The prescribed flat background has coordinates
$y_v^0(\sigma)$ and $h_e^0(\sigma)$ of positive $\sigma$ order.
The two endpoint sections have the same $h_e^0$, because their flat
restrictions use the common marked background, and all section
corrections vanished on the residual quotients.  As in
\cite[Section~6.5]{MT}, recenter by
\[
 y_v=y_v^0(\sigma)+z_v,\qquad h_e=h_e^0(\sigma)+z_e
\]
\emph{before} adjoining Laurent coefficients.  These substitutions
are continuous in the joint formal rings: every coefficient in $z$
is a $\sigma$-adically convergent sum.  Only afterward pass to
$K=\mathbb C((\sigma))$ and complete in the new internal coordinates.
Rename these coordinates $y_v,h_e$.

The shifted $F_e$ and $G_e$ remain algebraic germs over $K$, by
substitution into their finite algebraic presentations.  The
background is central-critical and its port holonomies are trivial,
so $G_e(0)=1$.  Also $F_e(0)=0$: along the original formal background
curve its derivative with respect to $\sigma$ is zero, since the
central gradient vanishes there, and its value at $\sigma=0$ is zero.
The full-column minors, the injective tangent minors, and the
determinant in \eqref{for:tangent-isomorphism} all had nonzero constant
term at the joint origin.  Evaluating at the background leaves units
in $\mathbb C[[\sigma]]$, hence nonzero elements of $K$.

We keep the notation $r_v,\mathcal I_v$ for the recentered list and
ideal in $T_v[[y_v]]$.  They may now have linear terms.  We shall not
again use their earlier quadratic order.  Subsequent substitutions
into Laurent-coefficient series are centered, or have only nilpotent
parameter constants; no such series is translated by a nonzero
$\sigma$-dependent constant after this passage.

\subsection{Four-slot laws in the residual ideals}

We now use the marked laws for the first time in the formal
calculation.  Fix $v$ and a law set $L\in\mathcal L_v$, ordered as
two red--blue pairs.  Let $X_1,\ldots,X_4$ be lifts of the based
port logarithms to $T_v[[y_v]]$ that vanish at the recentered
background.  Fix these four lifts once for the entire family of laws
on $L$.  Evaluate each marked word using those lifts and its prescribed
extra-loop conjugators.  This off-shell word expression agrees with
the actual law holonomy in the flat quotient; no off-shell homotopy
invariance of path holonomy is assumed.  Define $\mathcal J_L$ to be
the ideal generated by all
products of one matrix coordinate from each $X_j$.

Every corresponding law holds in the flat quotient by
$\mathcal I_v$, because its marked curve bounds a disk in $V_v$.
Before making that quotient, each entry of a law minus identity
belongs to $\mathcal J_L$.  Indeed the law is Brunnian in its four
slots: setting any one slot equal to zero makes the logarithmic
expression trivial.  Its leading terms, up to the fixed signs and an
invertible overall conjugation, are scalar entries of
\begin{equation}
\label{for:four-linear-law}
 [[\operatorname{Ad}(g_{a_1})X_1,\operatorname{Ad}(g_{a_2})X_2],
   [\operatorname{Ad}(g_{a_3})X_3,\operatorname{Ad}(g_{a_4})X_4]],
\end{equation}
with $a_1,a_2,a_3,a_4\in\{1,\ldots,9\}$ chosen independently.
Internal basing changes
in the marked body are words in its four tips; they are identity when
the tips vanish and change only higher terms.  Every higher term
belongs to $\mathfrak m_v\mathcal J_L$, where now
$\mathfrak m_v=(y_v,t)$ is the recentered maximal ideal.  These
statements are first identities in independent logarithms and then
remain ideal statements under their substitution in the vertex ring.

The multilinear tensor $[[U_1,U_2],[U_3,U_4]]$ on
$\mathfrak{sl}_2$ is nonzero.  For example, with
$[h_0,e_0]=2e_0$ and $[h_0,f_0]=-2f_0$, its value at
$(h_0,e_0,h_0,f_0)$ is $-4h_0$.  Because the prescribed adjoints
span the full endomorphism algebra over $K$, their independent
actions on the four inputs of this nonzero tensor span every scalar
coordinate product.  To see this directly, choose input vectors on
which the tensor is nonzero and an output functional detecting it;
rank-one endomorphisms send any selected input coordinate to those
vectors.  Expanding each such endomorphism in the prescribed adjoints
expresses the desired coordinate product as a linear combination of
the scalar tests in \eqref{for:four-linear-law}.

Let $Q_L$ be the ideal generated by the entries of the actual marked
law matrices minus identity, for all the specified conjugators and
signs; these are the relations, not merely their leading terms.
The preceding calculation gives
\[
 Q_L\subset\mathcal J_L\cap\mathcal I_v,
 \qquad \mathcal J_L=Q_L+\mathfrak m_v\mathcal J_L.
\]
Nakayama's lemma for the finitely generated module
$\mathcal J_L/Q_L$ yields
\begin{equation}
\label{for:four-product-ideal}
 \mathcal J_L\subset\mathcal I_v.
\end{equation}
This is the four-slot, mixed-endpoint version of
\cite[Lemma~6.6]{MT}; the proof uses neither reducedness nor flatness
over the parameter ring.

Partition $L=I\amalg J$ by the plus and minus designations.  In the
flat quotient, every plus logarithm is $-t_iT_i$.  Choose a nonzero
coordinate of each $T_i$.  At a minus end use instead the variable
port holonomy, which is its marked $\beta_j$ holonomy.  Since each
entry of $\exp(X_j)-1$ is in the ideal generated by the entries of
$X_j$, \eqref{for:four-product-ideal} implies that
\begin{equation}
\label{for:mixed-holonomy}
 \left(\prod_{i\in I}t_i\right)
 \prod_{j\in J}(G_j-1)_{a_jb_j}(M_v)
                    \in\mathcal I_v
\end{equation}
for every choice of one matrix entry at each minus end.  Here the
algebraic $G_j$ agrees with the indicated holonomy on the section
modulo $\mathcal I_v$, since that end uses the central slice.  Extra
nilpotents from plus ends outside $L$ remain in the same ring and do
not alter this argument.

\subsection{Polynomial perturbations and completion of the proof}

The remaining perturbation coefficients $\ell_e(G_e)$ contain formal
logarithms.  Replace them by the algebraic germs
\begin{equation}
\label{for:polynomial-weight}
 b_e=-\epsilon_e\operatorname{tr}\left(
 T_e\bigl((G_e-1)-\tfrac12(G_e-1)^2\bigr)\right).
\end{equation}
They vanish at the origin.  We must restore exact vanishing after
this replacement while retaining \eqref{for:mixed-holonomy}.

At a plus end, set its own parameter $t_e$ equal to zero while
retaining all other parameters of that vertex.  In the residual
quotient the edge is then central-critical and has both peripheral
holonomies trivial: $\alpha_e$ is prescribed trivial, and $\beta_e$
bounds in $K_v$.  The marking property of $G_e$ therefore gives
\[
 G_e(M_v)-1\in(\mathcal I_v,t_e).
\]
This is a quotient-ring identity.  The earlier coordinate change
linear in $t_e$ becomes identity under this specialization.
The logarithmic remainder $R_e=\ell_e(G_e)-b_e$ has order at least
three in the entries of $G_e-1$, and its internal derivatives have
order at least two there.  Consequently
\begin{equation}
\label{for:replacement-estimates}
 t_eR_e(M_v)\in t_e\mathcal I_v^3,
 \qquad t_e(\partial_hR_e)(M_v)\in t_e\mathcal I_v^2,
\end{equation}
because $t_e^2=0$.  Write $\mathfrak t_v=(t_e:e\in E_+(v))$.
The replaced function $\mathcal S_v$ has value in
$\mathfrak t_v\mathcal I_v^2$ on the existing section, and its
gradient is a matrix times the same fixed list $r_v$, still with a
unit full-column minor.

Apply the nilpotent correction of \cite[Lemma~3.9]{MT}.  Its mechanism
also explains why the recentered Hessian need not vanish.  An error
in $\mathfrak t_v^n\mathcal I_v^2$ is canceled to first order by a
displacement in $\mathfrak t_v^n\mathcal I_v$, using the unit minor.
The Taylor remainder belongs to
$\mathfrak t_v^{2n}\mathcal I_v^2$, and the gradient coefficient
matrix changes by a $\mathfrak t_v^n$-multiple.  Since
$\mathfrak t_v$ is nilpotent, iteration terminates.  The resulting
section still agrees with the old one modulo $\mathcal I_v$, and
\[
 \mathcal S_v(M_v)=0,
 \qquad \bigl((\partial_h\mathcal S_v)(M_v)\bigr)=\mathcal I_v.
\]
There is no correction at a vertex with no plus ends.  In every case
the central tangent is unchanged.  The corrections allow all products
of distinct parameters; they do not replace $\mathfrak t_v$ by a
square-zero ideal.

Each $b_j$ in \eqref{for:polynomial-weight} is a polynomial without
constant term in entries of $G_j-1$.  Equation
\eqref{for:mixed-holonomy} therefore gives
\eqref{for:mixed-membership}, and the last corrections preserve it
because they are identity modulo $\mathcal I_v$.

Finally, these sections are smooth formal graphs over the full ring
$T_v$.  Choose a linear projection of the incident edge coordinates
that is invertible on the central tangent of $M_v$.  Its composite
with $M_v$ has invertible Jacobian over $T_v[[y_v]]$, since it does
modulo $\mathfrak t_v$.  Its constant term may lie in the nilpotent
ideal $\mathfrak t_v$; translation by such a constant is a legitimate
formal substitution.  The formal inverse theorem thus puts the
section into graph form.  The same reparametrization transports all
memberships.  The central tangent difference remains the isomorphism
\eqref{for:tangent-isomorphism}, transported through common invertible
central edge coordinates.  This proves Proposition~\ref{for:formal-data}.

\begin{proof}[Proof of Proposition~\ref{for:odd-nonfilling}]
A marked filling would give the cut pieces of Lemma~\ref{for:cuts},
and hence the formal data of Proposition~\ref{for:formal-data} on
$\Gamma_d$.  Its law family consists of every choice of two red and
two blue incident edges.  Corollary~\ref{ten:odd-cycle} forbids exactly
these data when $d$ is odd.  This contradiction proves the proposition.
\end{proof}
